\section{Introduction}\label{sec:introduction}

The Milnor number measures the topology of a nearby smooth level of an
isolated hypersurface singularity: its Milnor fiber has the homotopy type
of a bouquet of spheres, with one sphere for each unit of the Milnor
number \cite[Theorems~6.5 and~7.2]{Milnor1968}; the analytic
formulation in terms of the Jacobian algebra is recalled in
\cite[Section~1, Theorem~1.5]{LeRamanujam1976}.  A basic question in
equisingularity theory asks whether keeping this number constant also
keeps the defining function topologically unchanged.  For complex
surfaces, this requires a method that works beyond the dimensions in
which the classical cobordism argument applies.

Let $\Delta\subset\mathbb C$ be a disk centered at zero and let
\[
 F:(\mathbb C^3\times\Delta,\{0\}\times\Delta)\longrightarrow(\mathbb C,0)
\]
be holomorphic.  Write $f_t(x)=F(x,t)$ and suppose that $f_t$ has an
isolated critical point at the origin.  Its Milnor number is
\[
 \mu(f_t,0)=\dim_{\mathbb C}
 \frac{\mathcal O_{\mathbb C^3,0}}
 {(\partial f_t/\partial x_1,\partial f_t/\partial x_2,
   \partial f_t/\partial x_3)}.
\]
The family is \emph{topologically right-trivial} if a homeomorphism germ
of the ambient family, preserving the parameter and the origin section,
identifies each defining function with the central one.  The coordinate
changes and their inverses must vary jointly continuously with the
parameter.

\begin{theorem}\label{thm:main}
Let $F:(\mathbb C^3\times\Delta,\{0\}\times\Delta)\to(\mathbb C,0)$
be holomorphic.  Suppose that, for every $t\in\Delta$, $f_t(0)=0$,
$f_t$ has an isolated critical point at $0$, and
$\mu(f_t,0)=\mu(f_0,0)$.  After shrinking the disk and the germ
representatives, there is a homeomorphism germ
\[
 \Phi(x,t)=(\phi_t(x),t),\qquad
 \phi_t(0)=0,\quad \phi_0=\mathrm{id},\qquad
 f_t\bigl(\phi_t(x)\bigr)=f_0(x).
\]
Both $\Phi$ and $\Phi^{-1}$ are jointly continuous on neighborhoods
of the origin section over the original disk.
\end{theorem}

Thus the surface case of the $\mu$-constant problem has a positive
answer.  The parameter dependence may be arbitrary and holomorphic, and
the surface singularities need not be log canonical.  The finite
parameter covers used in the proof are removed before constructing the
homeomorphism in the theorem.  The same ambient maps identify the zero
sets and the nearby nonzero levels, keeping the value of the function
fixed throughout the parameter disk.

\subsection*{The problem and its antecedents}

L\^e and Ramanujam proved that a $\mu$-constant family of isolated
hypersurface singularities has constant embedded topological type when
the hypersurface dimension is different from two
\cite[Theorem~2.1]{LeRamanujam1976}.  Their argument compares Milnor
fibers through a homologically trivial cobordism.  In higher dimensions
the $h$-cobordism theorem turns this comparison into a product; the
curve case follows from surface topology.  Complex dimension two is
the dimension in which this argument leaves the product structure
unresolved.  Outside complex dimension two, Timourian and King established local
triviality of families of function germs
\cite{Timourian1977,King1980}.  This additional passage matters here:
constancy of the individual topological types and a jointly continuous
choice of ambient coordinate changes are different conclusions.

Whitney regularity supplies a stronger geometric form of equisingularity.
For a hypersurface germ in three variables, let $\mu^{(j)}$ be the
Milnor number of its restriction to a general $j$-dimensional linear
subspace.  Then $\mu^{(3)}=\mu$, while $\mu^{(1)}$ is the multiplicity
minus one and $\mu^{(2)}$ records the generic plane section.  Teissier
proved that constancy of this full sequence implies Whitney regularity
along the singular section \cite[Chapter~II, Theorem~3.9]{Teissier1973};
Brian\c con and Speder proved the converse
\cite[Theorem~2]{BrianconSpeder1976}.  Their examples also show that
topological triviality need not imply Whitney regularity
\cite{BrianconSpeder1975}\cite[Section~3]{BrianconSpeder1976}.
Thus a proof under the sole $\mu$-constant hypothesis must account for
families whose generic plane-section Milnor number changes.

The equimultiplicity theorem of Fern\'andez de Bobadilla and Pe\l ka
settles a separate part of this numerical comparison: a family with
constant finite Milnor number has constant multiplicity in every
dimension \cite[Theorem~1.1]{FernandezBobadillaPelka2024}.  Their theorem
allows even formal power series whose coefficients vary continuously
along a real interval.  Its proof gives A'Campo's radius-zero
monodromy model a symplectic structure and extends McLean's
Floer-theoretic spectral sequence to compare multiplicities without
assuming a common Milnor radius.  For the present surface problem it supplies constancy of
$\mu^{(1)}$, but does not supply constancy of $\mu^{(2)}$.  We use it
with Teissier's criterion in the log-canonical case; the double-point
case reduces to the classical $\mu$-constant problem for plane curves.

Vector fields offer another approach to the parameter dependence.
The L\^e--Saito--Teissier criterion relates constancy of the Milnor
number to the smallness of the parameter derivative compared with the
spatial gradient \cite{LeSaito1973}\cite[Chapter~II, Remark~3.10]{Teissier1973}.
Such an estimate makes a gradient lift continuous at the singular
section, but continuity alone does not give unique, continuously
invertible transport.  Parusi\'nski constructed trivializing vector
fields for $\mu$-constant linear families $f(x)+tg(x)$ of isolated
hypersurface singularities, without a dimension restriction
\cite[Corollary~2.1]{Parusinski1999}.  Appendix~\ref{sec:continuous-gradient}
separates the gradient estimate from the trajectory control needed
for general parameter dependence.

For surfaces, simultaneous resolution gives a bridge from numerical
invariants to topology.  Wahl associated a logarithmic characteristic
number to a normal surface singularity \cite{Wahl1990}.
Okuma showed that, for a deformation of a non-log-canonical Gorenstein
surface singularity, constancy of this number produces simultaneous surface resolutions
after finite base change, with reduced exceptional divisors on the
fibers over the distinguished singular section
\cite[Theorem~1.2]{Okuma2004}.

Recent work describes further constraints on degenerations under the
$\mu$-constant hypothesis.  Aldasoro Rosales studies semistable models
after suitable base change and projective compactification.  Her
results include rational acyclicity of the special-fiber dual complex,
Tate type for the second cohomology of the additional vertical
surfaces, and equality of the first Betti numbers of the special fiber
and its distinguished component
\cite[Theorem~1.2]{AldasoroRosales2025}.  The comparison below likewise
studies these additional surfaces, using their boundary complements
and logarithmic canonical squares.  Peters proves symplectic
equivalence of completed Milnor fibers and contact equivalence of
links for $\mu$-constant polynomial families without vanishing folds
\cite[Theorem~3.7]{Peters2026}.  That extra hypothesis provides a common
Milnor radius.  Our argument instead obtains transport from a resolved
zero family and extends it to the ambient function levels.

\subsection*{The resolution invariant and the numerical comparison}

Write $X_t=\{f_t=0\}$.  These are normal Gorenstein surfaces.  For
non-log-canonical germs, our main numerical task is to deduce constancy
of Wahl's logarithmic surface invariant from $\mu$ alone.
On a good resolution, write the relative Zariski decomposition of the
canonical divisor plus the full reduced exceptional divisor as $P+N$,
with $P$ relatively nef and $N$ effective and exceptional.  The
exceptional intersection matrix identifies a unique exceptional rational
cycle $P_{\mathrm{loc}}$ having the same intersections as $P$ with every
exceptional curve.  The invariant is
\[
 \beta=-P_{\mathrm{loc}}^2.
\]
Section~\ref{sec:local-square} recalls its resolution independence and
the conventions used by Okuma \cite[Section~2]{Okuma2004}.  It also
expresses this local square as a difference of ordinary projective
squares.

The proof rests on a cancellation of Euler characteristics.  To explain
it, first consider a projective family with just the prescribed surface
singularity in each fiber.  A semistable log minimal model replaces its
special fiber by a component $S_0$ dominating $X_0$ and additional
projective surfaces $S_a$ over the singular point.  Each component $S$
has a reduced boundary $B_S$, consisting of its intersections with the
other vertical components and the horizontal exceptional divisors.
Write $S^\circ=S\setminus B_S$.  The construction gives
$S_0^\circ\simeq X_0\setminus\{0\}$ and makes every
$K_{S_a}+B_{S_a}$ nef.

Constancy of $\mu$ keeps the Euler characteristic of the punctured
projective fibers constant.  Their nearby Euler comparison, after the
contribution of $S_0^\circ$ cancels, is
\[
 0=\sum_a\chi_c(S_a^\circ)+\delta,\qquad \delta\ge0.
\]
Here $\chi_c$ denotes Euler characteristic with compact supports, and
$\delta$ is the sum of the excess Euler characteristics of the local
smoothings at interior singular points of the extra surfaces.
Greuel--Steenbrink's first-Betti-number theorem makes these excesses
nonnegative \cite{GreuelSteenbrink1983}.
Langer's orbifold Bogomolov--Miyaoka--Yau inequality supplies the other
sign \cite{Langer2003}:
\[
 0\le (K_{S_a}+B_{S_a})^2\le3\chi_c(S_a^\circ)
 \qquad\text{for every }a.
\]
Thus every extra logarithmic square vanishes.  Specialization of
projective intersections then gives constancy of $\beta$.  The key
point is to retain a boundary for which both the Euler comparison and
the surface inequality hold.

The technical development establishes these properties in order.
Section~\ref{sec:algebraic-reduction} reduces this scalar assertion to
polynomial families over algebraic curves.  Individual finite
determinacy supplies polynomial representatives, and constructibility
of the resolution invariant gives the curve comparison.  These
individual equivalences are used only to compare numerical invariants.
The high-degree compactification, following the construction in
\cite[Lemma~2.5]{Steenbrink1985}, gives projective surfaces with precisely
the prescribed singularity and a smooth complement.

Sections~\ref{sec:log-model}--\ref{sec:full-boundary} use semistable
reduction and the threefold log minimal model program
\cite[Chapter~7]{KM1998} to obtain a model with a reduced special
fiber.  The boundary retains every surviving horizontal exceptional
divisor.  Relative nefness and the negativity lemma identify its full
support over the singular section.  Adjunction then supplies precisely
the reduced boundaries used in the comparison above.

Sections~\ref{sec:nearby-euler}--\ref{sec:surface-estimates} establish
the Euler comparison and its signs.  The essential local assertion is
that a boundary intersection contributes zero to the nearby Euler
function, including at singular points of the threefold.  Finite covers
that make the boundary Cartier and the structure of threefold dlt strata reduce
this assertion to the Euler calculation for a product of coordinates.
Proper compatibility of nearby cycles then permits the local
calculation to be integrated over the special fiber
\cite{SchuermannNearby,MaximSchuermann}.  The interior correction and the
comparison of orbifold and ordinary Euler characteristics complete the
two sign calculations.  Section~\ref{sec:square-specialization} converts
the resulting vanishing into constancy of $\beta$.

The log-canonical branch is handled separately.  Constancy of the
spectrum, following Varchenko and Steenbrink
\cite{Varchenko1982,Steenbrink1985}, together with the threshold formula
\cite[Theorem~9.5]{Kollar1997Pairs} and inversion of adjunction
\cite{Kawakita2007}, shows that log canonicity is constant along the
family.  Section~\ref{sec:low-cases} treats this branch and the
double-point case through sectional Milnor numbers.

\subsection*{From resolution to a homeomorphism of the functions}

In the non-log-canonical branch, Okuma's theorem produces a simultaneous
minimal semigood resolution
after a finite disk cover \cite[Theorem~1.2]{Okuma2004}.  Its reduced
exceptional restrictions are essential: they allow the remaining
relative-normal-crossings assertion to be proved by the Euler
characteristic of the exceptional curves.  A node smoothing would
lower that characteristic, whereas the resolved tubes and the
punctured zero fibers make it constant.

Section~\ref{sec:ambient-lifting} constructs complete lifts of the
parameter directions on the resolved tubes.  These descend to the
punctured zero fibers and extend to the ambient nonzero levels while
preserving the defining function.  Properness of the tubes and reverse
transport establish joint continuity at the origin for both the maps
and their inverses.  An explicit correction along the angular seam of
$t=u^e$ then yields a trivialization over the original disk.
Section~\ref{sec:low-cases} supplies the double-point and log-canonical
cases using sectional Milnor numbers and completes the proof.


Appendix~\ref{app:local-numerical} supplies the marked monomial
construction used in Section~\ref{sec:log-model}.  It also gives two
variants of the numerical argument: specialization of intersection
squares supported properly over a curve and comparison within an
irreducible component of a finite jet space.
Appendix~\ref{sec:continuous-gradient} gives an ambient transport
criterion using continuous gradient lifts and explicit control of
trajectories at the singular section.  Its two sufficient geometric
conditions come from Whitney regularity and from a proper simultaneous
resolution with relative normal crossings.  The criterion separates the
smallness of the vector fields from the additional control needed to
obtain continuous flows and continuous inverse flows.
